\chapter{Resume executif}

Ce rapport presente le developpement de \textit{Sport Academy Pro}, une plateforme multi-sport de gestion des academies sportifs. L'ecosysteme logiciel est compose de quatre projets interconnectes:
\begin{itemize}
  \item un backend Java Spring Boot pour la gestion metier de l'academie;
  \item un frontend Flutter pour mobile et web admin;
  \item un chatbot Django base sur une recherche de reponses par similarite;
  \item un service IA FastAPI pour l'analyse du potentiel sportif et le scouting.
\end{itemize}

Le contenu est construit a partir des documents et fichiers techniques existants dans les repertoires \texttt{football-academy}, \texttt{moez\_project}, \texttt{chatbot} et \texttt{scouting-ai-service}.

\section{Contexte et problematique}

Les academies et clubs sportifs multi-sports font face a des defis majeurs: donnees dispersées, faible traçabilité, organisation sportive complexe, suivi financier insuffisant, communication non structurée et besoin d'outils décisionnels. \textit{Sport Academy Pro} vise à centraliser et digitaliser ces processus via une solution multi-sport extensible.

\section{Objectifs du travail}

Les objectifs couverts par ce rapport sont les suivants:
\begin{itemize}
  \item documenter l'architecture globale et les interactions entre services;
  \item decrire les choix techniques backend, frontend, chatbot et service IA;
  \item presenter les mecanismes de securite, de communication temps reel et de deploiement;
  \item detailler les fonctionnalites multi-sport et les capacites de scouting IA;
  \item synthese des limitations connues et des pistes d'amelioration.
\end{itemize}

\section{Resultats techniques principaux}

\begin{itemize}
  \item \textbf{Backend}: Spring Boot 3.2.4, Java 17, securise par JWT, avec endpoints REST et WebSocket STOMP pour le chat temps reel.
  \item \textbf{Frontend}: Flutter avec deux points d'entree: application mobile utilisateur et plateforme web admin, supportant trois langues (FR, EN, AR).
  \item \textbf{Chat}: Chat interne optimise (pattern singleton, pre-subscriptions, mise a jour temps reel) avec messagerie asynchrone via JMS.
  \item \textbf{Chatbot}: Service Django expose via \texttt{POST /api/chat}, avec pipeline exact match, fuzzy match et TF-IDF pour l'assistance utilisateurs.
  \item \textbf{Service IA}: Microservice FastAPI fournissant ML01 (score de potentiel), ML02 (analyse evolution), ML03 (prediction churn) et SC01 (scouting avance).
  \item \textbf{Multi-sport}: Architecture configurable supportant football, basketball, handball, etc., avec positions et statistiques adaptables.
  \item \textbf{DevOps}: Base de deploiement conteneurisee via Docker Compose, pipeline CI/CD GitLab, et monitoring des services.
\end{itemize}

\section{Fonctionnalites metier couvertes}

\begin{itemize}
  \item \textbf{F01-F13}: Authentification securisee, gestion profils, equipes, calendrier, statistiques, paiements, notifications, messagerie, chatbot, recherche, back-office et integration n8n.
  \item \textbf{ML01-ML03}: Analyse du potentiel sportif avec explications des facteurs, analyse de l'evolution (progression/stabilite/regression), et prediction du risque d'abandon avec actions recommandees.
  \item \textbf{MS01-MS03}: Activation et configuration multi-sport avec UX unifiee.
  \item \textbf{SC01}: Recherche avancee de joueurs, comparaison selon statistiques et tendances, et generation de shortlists professionnelles.
\end{itemize}

\section{Valeur du rapport}

Ce projet LaTeX fournit une base operationnelle pour votre rapport PFE: vous pouvez le compiler immediatement sur Overleaf ou en local, puis remplacer progressivement les elements de personnalisation (noms, captures d'ecran, diagrammes, resultats de tests). La structure couvre integralement les exigences du cahier des charges du Master DSIR.
